
\begin{table*}[tp]
	\centering
	\setlength{\abovedisplayskip}{0pt}
	\setlength{\belowdisplayskip}{0pt}
    \small
	\caption{\textbf{Quantitative comparison of novel view synthesis on RE10K~\cite{zhou2018stereo}.} Our method achieves comparable performance with previous SOTA methods on few-view settings, and outperforms them with more views input. }
    \vspace{2pt}
    \scalebox{0.9}{
	\begin{tabular}{l|w{c}{1.0cm}w{c}{1.0cm}w{c}{1.1cm}w{c}{1.0cm}w{c}{1.0cm}w{c}{1.1cm}w{c}{1.0cm}w{c}{1.0cm}w{c}{1.1cm}}
		\toprule
		 \multirow{2}{*}{Method} & \multicolumn{3}{c}{2 views} & \multicolumn{3}{c}{4 views} & \multicolumn{3}{c}{8 views}  \\
          & PSNR $\uparrow$ & SSIM $\uparrow$ & LPIPS $\downarrow$ & PSNR $\uparrow$ & SSIM $\uparrow$ & LPIPS $\downarrow$ & PSNR $\uparrow$ & SSIM $\uparrow$ & LPIPS $\downarrow$ \\
		\midrule

        MVSplat~\cite{chen2024mvsplat}&24.73&0.829&0.171&21.91&0.753&0.358&20.41&0.696&0.434\\
        NoPoSplat~\cite{ye2024no}&\underline{25.59}&\textbf{0.849}&\textbf{0.140}&22.63&0.771&0.302&21.08&0.671&0.412 \\
        FLARE~\cite{zhang2025flare}&24.83&0.834&0.165&\underline{23.16}&\underline{0.794}&\underline{0.218}&22.41&0.765&0.314 \\
        \midrule
        Spann3R~\cite{wang20243d}+GS&22.43&0.761&0.317&22.07&0.732&0.374&21.86&0.701&0.396 \\
        CUT3R~\cite{wang2025continuous}+GS&23.12&0.798&0.208&22.91&0.775&0.241&\underline{22.78}&\underline{0.774}&\underline{0.236} \\
        \textbf{Ours}&\textbf{25.78}&\underline{0.845}&\underline{0.144}&\textbf{25.56}&\textbf{0.839}&\textbf{0.147}&\textbf{25.59}&\textbf{0.841}&\textbf{0.148} \\

		\bottomrule
	\end{tabular}}
    \label{tab:re10k}
\end{table*}

\begin{table*}[tp]
	\centering
	\setlength{\abovedisplayskip}{0pt}
	\setlength{\belowdisplayskip}{0pt}
    \small
	\caption{\textbf{Quantitative comparison of novel view synthesis on ScanNet~\cite{dai2017scannet}.} Our method outperforms previous SOTA methods on all input view settings, even larger gains on 30-view setting.}
    \vspace{2pt}
    \scalebox{0.9}{
	\begin{tabular}{l|w{c}{1.0cm}w{c}{1.0cm}w{c}{1.1cm}w{c}{1.0cm}w{c}{1.0cm}w{c}{1.1cm}w{c}{1.0cm}w{c}{1.0cm}w{c}{1.1cm}}
		\toprule
		 \multirow{2}{*}{Method} & \multicolumn{3}{c}{10 views} & \multicolumn{3}{c}{20 views} & \multicolumn{3}{c}{30 views}  \\
          & PSNR $\uparrow$ & SSIM $\uparrow$ & LPIPS $\downarrow$ & PSNR $\uparrow$ & SSIM $\uparrow$ & LPIPS $\downarrow$ & PSNR $\uparrow$ & SSIM $\uparrow$ & LPIPS $\downarrow$ \\
		\midrule
                
        MVSplat~\cite{chen2024mvsplat}&20.03&0.686&0.414&18.72&0.561&0.511&16.02&0.505&0.591\\
        NoPoSplat~\cite{ye2024no}&21.34&0.729&0.360&19.42&0.635&0.469&18.73&0.574&0.502 \\
        FLARE~\cite{zhang2025flare}&21.75&0.751&0.330&19.91&0.652&0.447&19.15&0.619&0.485 \\
        \midrule
        Spann3R~\cite{wang20243d}+GS&21.54&0.732&0.357&\underline{20.91}&\underline{0.712}&\underline{0.367}&\underline{20.42}&\underline{0.689}&\underline{0.387} \\
        CUT3R~\cite{wang2025continuous}+GS&\underline{21.97}&\underline{0.762}&\underline{0.323}&20.79&0.709&0.382&20.01&0.687&0.401 \\
        \textbf{Ours}&\textbf{24.13}&\textbf{0.811}&\textbf{0.174}&\textbf{23.86}&\textbf{0.782}&\textbf{0.189}&\textbf{23.73}&\textbf{0.769}&\textbf{0.194} \\
		\bottomrule
	\end{tabular}}
    \label{tab:scannet}
\end{table*}



\section{Experiments}
\subsection{Experimental Settings}
\label{sec:exp}
\paragraph{Datasets.} We train and evaluate our model on two widely-used real-world datasets: RealEstate10k (RE10K)~\cite{zhou2018stereo} and ScanNet~\cite{dai2017scannet}. As our primary benchmarks, RE10K is used for performance evaluation on video sequences with a limited spatial scope, while ScanNet serves as the basis for our room-scale reconstruction task. To evaluate the generalization ability of our model, we further perform zero-shot evaluation on the DL3DV dataset~\cite{ling2024dl3dv}.

\paragraph{Baselines.} To assess the effectiveness of our OnlineX framework, we compare it against several state-of-the-art feed-forward 3D reconstruction methods, which fall into two main categories: (1) offline feed-forward 3DGS approaches, including MVSplat~\cite{chen2024mvsplat}, NoPoSplat~\cite{ye2024no} and FLARE~\cite{zhang2025flare}; and (2) online feed-forward pointmap prediction methods, such as Spann3R~\cite{wang20243d} and CUT3R~\cite{wang2025continuous}. For the offline 3DGS approaches, we provide all input views simultaneously, allowing them to perform an offline reconstruction. This setting is inherently less challenging than our online scenario, which must build the scene progressively from a sequential stream without access to future frames. For the online pointmap methods, since they do not natively produce Gaussian representations, we adapt them for a fair comparison by extending their architectures with a Gaussian Splatting prediction head identical to our own and subsequently fine-tuning them on the corresponding datasets. In addition, for the 3D scene understanding task, we compare our method to the per-scene optimization-based method LangSplat~\cite{qin2024langsplat} and Gaussian-Grouping (GS-Group)~\cite{ye2024gaussian} with the same input views.  



\paragraph{Implementation Details.} During training, we sample sequences of randomly varying lengths (from 4 to 15 views) to teach the model the principles of online, iterative modeling, thereby ensuring its ability to generalize to longer, unseen sequences at inference time. he sampling interval between neighboring frames is set to 10 for RE10K and 20 for ScanNet, resulting in a moderate view overlap. We train our OnlineX model using the AdamW optimizer~\cite{loshchilov2017fixing} with an initial learning rate of $5 \times 10^{-5}$ for a total of 30,000 iterations on 4 NVIDIA RTX A6000 GPUs with an effective batch size of 8. The loss weights $\lambda_{\text{aux}}$, $\lambda_{1}$, $\lambda_{2}$ and $\lambda_{3}$ are set to 0.8, 1, 1 and 0.5, respectively.  We adopt pre-trained weights of MASt3R~\cite{leroy2024grounding} with adaptive adjustment. 



% \begin{wraptable}{r}{0.55\textwidth}
%     \centering
%     \caption{2D Quantitative Results on ScanNet Dataset. We
%     report the open-vocabulary 2D semantic segmentation (IoU
%     scores). The bold denotes the best results.}
%     \vspace{2pt}
%     \begin{tabular}{c|c|c|c}
%         \toprule
%         Scene & LangSplat~\cite{qin2024langsplat} & GS-Group~\cite{ye2024gaussian} & Ours \\
%         \midrule
%         085 & 6.19 & 18.92 & \textbf{52.70} \\
%         114 & 13.19 & 13.35 & \textbf{30.62} \\
%         616 & 08.12 & 09.47 & \textbf{36.44} \\
%         617 & 11.39 & 10.62 & \textbf{33.05} \\
%         \midrule
%         Mean & 9.72 & 13.09 & \textbf{38.20} \\
%         \bottomrule
%     \end{tabular}
%     \label{tab:scannet_lang}
% \end{wraptable}


% \begin{wraptable}{r}{0.55\textwidth}
%     \centering
%     \caption{Runtime comparison of LangSplat, GS-Group and Ours. The bold denotes the best results.}
%     \vspace{2pt}
%     \begin{tabular}{c|c|c|c}
%         \toprule
%         Method & LangSplat~\cite{qin2024langsplat} & GS-Group~\cite{ye2024gaussian} & Ours \\
%         \midrule
%         Runtime & 25min6s & 10min10s & \textbf{10s} \\
%         \bottomrule
%     \end{tabular}
%     \label{tab:scannet_lang_time}
% \end{wraptable}





\begin{figure*}[t]
    \centering
    \includegraphics[width=1.0\linewidth]{figures/vis_1.pdf}
    \caption{\textbf{Qualitative comparison for novel view synthesis on RE10K (top two rows) and ScanNet (bottom two rows).} We adopt the 4-view setting for RE10K and 15-view setting for ScanNet. }
    \label{fig:vis}
\end{figure*}

\begin{table}[t]
	\centering
	\small
	\caption{\textbf{Camera pose estimation on ScanNet.} We report Absolute Translation Error (ATE), Relative Translation Error (RPE trans), and Relative Rotation Error (RPE rot).}
	\label{tab:pose_scannet}
	\begin{tabular}{l ccc}
		\toprule
		Method & ATE$\downarrow$ & RPE trans$\downarrow$ & RPE rot $\downarrow$ \\
		\midrule
		Spann3R~\cite{wang20243d} & \underline{0.096} & 0.023 & 0.661 \\
		CUT3R~\cite{wang2025continuous} & 0.099 & \underline{0.022} & \underline{0.600} \\
		\textbf{Ours} & \textbf{0.085} & \textbf{0.019} & \textbf{0.550} \\
		\bottomrule
	\end{tabular}
\end{table}
\begin{table}[t]
	\centering
	\setlength{\abovedisplayskip}{0pt}
	\setlength{\belowdisplayskip}{0pt}
    \small
	\caption{\textbf{Quantitative comparison of semantic segmentation on ScanNet~\cite{dai2017scannet}.} We report the average IoU scores ($\%$) and average accuracy ($\%$). We denote Gaussian-Grouping~\cite{ye2024gaussian} as GS-Group.}
    \vspace{2pt}
    \scalebox{0.9}{
    \begin{tabular}{l|cccc}
		\toprule
		 \multirow{2}{*}{Method} & \multicolumn{2}{c}{5 views} & \multicolumn{2}{c}{15 views}  \\
          & mIoU $\uparrow$ & mAcc $\uparrow$ & mIoU $\uparrow$ & mAcc $\uparrow$ \\
		\midrule
        
        LangSplat~\cite{qin2024langsplat}&\underline{54.63}&\underline{71.15}&\underline{56.91}&\underline{73.41}\\
        GS-Group~\cite{ye2024gaussian}&51.61&68.32&53.26&67.63  \\
        \textbf{Ours}&\textbf{58.83}&\textbf{77.12}&\textbf{57.79}&\textbf{76.36} \\
		\bottomrule
	\end{tabular}}
    \label{tab:scannet_lang}
    \vspace{-2pt}
\end{table}


\subsection{Results}
\paragraph{Novel View Synthesis.}
Table \ref{tab:re10k} and Table \ref{tab:scannet} present the novel view synthesis (NVS) results of our proposed OnlineX on the RE10K~\cite{zhou2018stereo} and ScanNet~\cite{dai2017scannet} datasets, respectively. For a comprehensive evaluation, we test on varying sequence lengths: 2, 4, and 8 views for RE10K, and 10, 20, and 30 views for ScanNet. In few-view settings, our method achieves performance comparable to offline reconstruction baselines. As the number of views increases, OnlineX demonstrates stable performance and consistently outperforms competing online approaches. These results demonstrate our model's strong performance in both sparse-view and long-term online scenarios.


\begin{figure*}[t]
    \centering
    \includegraphics[width=1.0\linewidth]{figures/vis_lang.pdf}
    \caption{\textbf{Qualitative comparison for semantic segmentation on ScanNet. } Here we showcase one scene with 15 input views. The masks predicted by ours contain more complete regions than other methods, such as the "Wall" prompt, which also surpasses the GT masks.}
    \label{fig:vis_lang}
    \vspace{3pt}
\end{figure*}

\begin{figure}[t]
    \centering
    \includegraphics[width=1.0\linewidth]{figures/zeroshot_1.pdf}
    \caption{\textbf{Qualitative results for zero-shot generalization on DL3DV. } Our model can easily transfer to out-of-distribution data.}
    \label{fig:ood}
    \vspace{-3pt}
\end{figure}

Figure \ref{fig:vis} further shows the visualization results of NVS compared with baseline methods. It can be observed that our method significantly surpasses the baselines in predicting accurate global geometry, capturing fine-grained details, and reducing artifacts caused by overlapped Gaussians. 

\paragraph{Camera Pose Estimation.}
We evaluate the camera pose estimation accuracy of our method on the ScanNet dataset~\cite{dai2017scannet} with 30 input views. Following standard protocols~\cite{chen2024leap, zhao2022particlesfm, zhang2024monst3r}, we report Absolute Translation Error (ATE), Relative Translation Error (RPE trans), and Relative Rotation Error (RPE rot) after Sim(3) alignment with the ground truth. We compare our approach against online methods that similarly do not require camera calibration, specifically CUT3R~\cite{wang2025continuous} and Spann3R~\cite{wang20243d}. As shown in Table~\ref{tab:pose_scannet}, our OnlineX framework consistently outperforms both baselines across all three metrics, demonstrating the effectiveness of our decoupled architecture in maintaining a more accurate and robust camera trajectory.

\paragraph{Open-Vocabulary Semantic Segmentation.} For open-vocabulary segmentation, we query the rendered 2D language feature map by computing the per-pixel cosine similarity between its features and a given text embedding. This process yields a confidence map, where regions with high similarity scores are taken as the final segmentation result for the queried object. We evaluate our method using mIoU and mAcc on the ScanNet datasets, as shown in Table \ref{tab:scannet_lang}. Compared with existing 3D language field methods such as LangSplat~\cite{qin2024langsplat} and Gaussian Grouping~\cite{ye2024gaussian}, our approach achieves superior performance across both metrics. Visualization results in Figure \ref{fig:vis_lang} further show that our method accurately segments objects with detailed boundaries, highlighting the strength of our unified visual and semantic understanding in capturing visual guidance. 



\paragraph{Cross-Dataset Generalization.} We also assess the zero-shot generalization ability of our model, where we train on RE10K~\cite{zhou2018stereo} and directly evaluate it on DL3DV~\cite{ling2024dl3dv} for NVS task with 6 input views. As shown in Figure \ref{fig:ood}, our method in general shows superior zero-shot performance on out-of-distribution data. This improved generalization is largely attributed to the unified online architecture in our method, which enhances its adaptability across diverse scene types and varying lengths of input sequences.

% \paragraph{Long Sequence Reconstruction.}

% We additionally show quantative and qualitative results of our whole scene reconstruction (~ 40 images) in Table \ref{tab:long_sequence} and Figure \ref{fig:long_sequence} compared with FLARE~\cite{zhang2025flare}. Our method achieves superior performance with FLARE on the ScanNet dataset with accurate global geometry and fine-grained textures.
% \begin{wraptable}{r}{0.5\textwidth}
%     \centering
%     \footnotesize
%     \vspace{-25pt}
%     \caption{\textbf{Quantitative comparison for inference runtime and FPS. } }
%     \vspace{2pt}
%     \begin{tabular}{l|w{c}{1.0cm}w{c}{1.0cm}}
%       \toprule
%       Method &  Time $\downarrow$ & FPS $\uparrow$  \\
%       \midrule
%       pixelSplat~\cite{charatan2024pixelsplat} & 12 & 14 \\ 
%       MVSplat~\cite{chen2024mvsplat} & 11 & 14 \\
%       \midrule
%       Splatt3R~\cite{smart2024splatt3r} & 10 & 14 \\
%       NoPoSplat~\cite{ye2024no} & 9 & 14 \\
%       FLARE~\cite{zhang2025flare} & 8 & 14 \\
%       Ours& 14 & 14 \\
%       \bottomrule
%     \end{tabular}
%     \label{tab:runtime}
%     \vspace{-10pt}
% \end{wraptable}
\paragraph{Runtime Analysis.}
Our method achieves 23 frames per second (FPS) on 256×256 inputs using a single NVIDIA RTX A6000 GPU, supporting real-time applications. The inference time and GPU memory usage are comparable to that of CUT3R~\cite{wang2025continuous} and are significantly faster and lower than Spann3R~\cite{wang20243d} as shown in Table \ref{tab:runtime}.

\begin{table}[t]
	\centering
    \small
	\caption{\textbf{Runtime and memory analysis.} We report the FPS and memory usage metric compared with other online methods.}
    \begin{tabular}{l|ccc}
		\toprule
		  Method & Spann3R~\cite{wang20243d} & CUT3R~\cite{wang2025continuous} & \textbf{Ours} \\
        \midrule
        FPS & 13.35  & \textbf{25.78} & \underline{23.12} \\
        Memory(GB)&32.73&\textbf{19.76}&\underline{21.64}\\
		\bottomrule
	\end{tabular}
    \label{tab:runtime}
\end{table}

\begin{table}[t]
    \centering
    \small
    \caption{\textbf{Quantitative results of the ablation study. } We report novel view synthesis metrics on ScanNet with 10 views.}
    \begin{tabular}{l|ccc}
		\toprule
		Method & PSNR $\uparrow$ & SSIM $\uparrow$ & LPIPS $\downarrow$\\
		\midrule
		w/o Relative Extractor & 20.85 & 0.712 & 0.368 \\
		w/o Anchor State & 19.92 & 0.685 & 0.441  \\
		w/o Implicit Transform & 16.52 & 0.533 & 0.552\\
		w/o Implicit GS Fusion & \underline{22.61} & \underline{0.778} & \underline{0.284} \\
		\midrule
		\textbf{Ours (Full Model)} & \textbf{24.13} & \textbf{0.811} & \textbf{0.174}  \\
      \bottomrule
    \end{tabular}
    \label{tab:ablation}
\end{table}

\subsection{Ablation Studies}

% As shown in Table \ref{tab:ablation} and Figure \ref{fig:ablation}, we conduct extensive ablation studies to verify the effectiveness of each component in our method, which include the relative token extractor (RTE), the global token projector (GTP), the adaptive gaussian fusion module (AGF), the joint modeling of poses and language features. We separately remove the relative network and the global network, both of which result in corrupted global structures and inaccurate viewpoint estimations. This demonstrates that the relative network provides relatively accurate initialization, while the global network is essential for effective global modeling. It can be also observed that the AGF module is crucial for reducing the redundancy of Gaussians and improving the quality of overlapped regions. We further remove the pose and language tokens along with their corresponding losses. Removing pose tokens results in regional misalignment, whereas removing language tokens leads to misalignment particularly around object boundaries. This demonstrates that pose tokens act as a critical bridge from relative to global representations, and that visual-language joint modeling provides complementary information to enhance performance.

We conduct extensive ablation studies to verify the effectiveness of our key components (Table \ref{tab:ablation} and Figure \ref{fig:ablation}). Removing the Relative Geometry Extractor leads to a loss of fine-grained detail and inaccurate poses. Discarding the Anchor State results in severe camera drift, confirming its necessity for long-term consistency. Replacing our Implicit Pose Transformation with an explicit one causes visible seams between frames, demonstrating its importance for seamless reconstruction. Finally, omitting the Implicit GS Fusion module produces blurrier results with noticeable overlapping Gaussians at object boundaries.



\begin{figure}[t]
    \centering
    \includegraphics[width=1.0\linewidth]{figures/ablation.pdf}
    \caption{\textbf{Qualitative results of the ablation study.} From left to right, we visualize the results of our full model and four variants: without the Relative Geometry Extractor, without the Anchor State, without the Implicit Pose Transformation, and without the Implicit GS Fusion.}
    \label{fig:ablation}
    \vspace{-5pt}
\end{figure}



